\def \Subject {تفسیرپذیری}
\def \Session { یک}
% \setcounter{chapter}{\Session-1}
\renewcommand{\chaptername}{سوال}
\chapter{\Subject}

\section{توصیف‌پذیری مدل‌های هوش مصنوعی با استفاده از \lr{Grad-CAM}}

\subsection{\lr{Grad-CAM}}

\subsubsection{پرسش و پاسخ}

\begin{enumerate}
	
	\item 
	زیرا \lr{Grad-CAM} همان ایده و هدف \lr{CAM} را دارد: تولید نقشه‌ی حرارتی\LTRfootnote{heatmap} که نشان دهد هر ناحیه از تصویر چقدر در تصمیم‌گیری مدل مؤثر است. 
	
	\lr{CAM} فقط در مدل‌هایی قابل استفاده است که لایه‌ی \lr{Global Average Pooling} به همراه لایه‌ی آخر \lr{Fully Connected} داشته باشند. 
	
	اما \lr{Grad-CAM} با بهره‌گیری از این ایده که به جای وزن‌های لایه‌ی \lr{FC} می‌توان از مشتق نمره‌ی کلاس نسبت به نقشه‌های ویژگی\LTRfootnote{feature maps} استفاده کرد، روی تقریباً هر مدل کانولوشنی بدون نیاز به تغییر معماری قابل اعمال است. 
	
	در واقع اگر مدلی دارای معماری مورد نیاز \lr{CAM} باشد، \lr{Grad-CAM} به فرم \lr{CAM} فرو می‌ریزد و خروجی آن‌ها بسیار مشابه خواهد بود.
	
	\item 
	هدف از اعمال \lr{ReLU} در انتهای \lr{Grad-CAM} این است که فقط نواحی‌ای در نقشه‌ی حرارتی نمایش داده شوند که احتمال کلاس مورد نظر را \textbf{تقویت} می‌کنند. 
	
	اگر گرادیان\LTRfootnote{gradient} یک ناحیه منفی باشد، یعنی افزایش فعالیت آن ناحیه باعث کاهش نمره‌ی کلاس هدف می‌شود؛ بنابراین آن ناحیه در جهت تصمیم فعلی مدل عمل نمی‌کند. چنین نواحی‌ای معمولاً برای تصمیم‌گیری فعلی مفید نیستند. 
	
	بنابراین با اعمال \lr{ReLU}، تأثیر مقادیر منفی حذف شده و نقشه‌ی نهایی فقط نواحی «مثبت» و مؤثر را نشان می‌دهد. 
	
	(دقت شود که گرادیان منفی لزوماً به معنای «تشخیص کلاس دیگر» نیست، بلکه یعنی آن ناحیه با تصمیم فعلی هم‌جهت نیست.)
	
	\item 
	\lr{Grad-CAM} نشان می‌دهد مدل برای پیش‌بینی یک کلاس به کدام نواحی تصویر توجه می‌کند. 
	
	اگر مدل به جای ویژگی‌های معنادار (مانند لباس یا ابزار کار در تشخیص «پرستار» در مقابل «دکتر»)، روی ویژگی‌های نادرست و جانبی مانند جنسیت، نژاد، پس‌زمینه یا \lr{watermark} تمرکز کند، این نواحی در نقشه‌ی حرارتی پررنگ دیده می‌شوند. 
	
	بنابراین با تحلیل نقشه‌های حرارتی روی نمونه‌های مختلف (به‌خصوص نمونه‌هایی که برچسب آن‌ها با ویژگی‌های زمینه‌ای در تضاد است) می‌توان وجود سوگیری\LTRfootnote{bias} را کشف کرد. 
	
	هرچند \lr{Grad-CAM} به تنهایی سوگیری را کمّی نمی‌کند، اما ابزار قدرتمندی برای بازشناسی بصری آن است.
	
	\item 
	\lr{Guided Grad-CAM} ترکیبی از دو روش است:
	
	\begin{itemize}
		
		\item 
		\textbf{\lr{Guided Backpropagation}}:
		در هنگام محاسبه‌ی گرادیان، در لایه‌های \lr{ReLU} علاوه بر شرط مثبت بودن ورودی، شرط مثبت بودن گرادیان برگشتی نیز اعمال می‌شود. این کار باعث می‌شود گرادیان نهایی تا سطح پیکسل‌های ورودی فقط از مسیرهایی عبور کند که هم فعال بوده‌اند و هم تأثیر مثبت داشته‌اند. 
		
		نتیجه، یک نگاشت با جزئیات پیکسلی بالا است که بیشتر ساختارهای محلی و لبه‌های تصویر را برجسته می‌کند؛ هرچند الزاماً تفکیک‌کننده‌ی کلاس نیست.
		
		\item 
		\textbf{\lr{Grad-CAM}}:
		نقشه‌ی حرارتی با رزولوشن پایین که نواحی مهم را مشخص می‌کند، اما به دلیل بزرگ‌نمایی\LTRfootnote{upscale} شدن، فاقد دقت پیکسلی است.
		
	\end{itemize}
	
	در \lr{Guided Grad-CAM}، خروجی \lr{Guided Backpropagation} به‌صورت نقطه‌به‌نقطه در خروجی \lr{Grad-CAM} (که به همان ابعاد بزرگ‌نمایی شده است) ضرب می‌شود. 
	
	نتیجه، نقشه‌ای با \textbf{جزئیات پیکسلی بالا} و در عین حال \textbf{تفکیک‌کننده‌ی کلاس}\LTRfootnote{class-discriminative} است. 
	
	به بیان ساده‌تر، از \lr{Grad-CAM} به‌عنوان \textbf{ماسک}\LTRfootnote{mask} استفاده می‌شود که تعیین می‌کند «کدام نواحی مهم هستند» و از \lr{Guided Backpropagation} برای مشخص کردن «کدام پیکسل‌ها» در آن نواحی اهمیت دارند استفاده می‌شود.
	
\end{enumerate}

\subsubsection{پیاده‌سازی}

\begin{figure}[H]
	
	\centering
	\includegraphics[width=\textwidth]{images/q1_samples.png}
	\caption{
		پنج نمونه از مواردی که مدل به‌درستی کلاس آن‌ها را تشخیص داده است.
	}
	\label{q1:samples}
	
\end{figure}

\begin{figure}[H]
	
	\centering
	\includegraphics[width=\textwidth]{images/q1_gradcam.png}
	\caption{
		خروجی 
		\lr{Grad-CAM}
		برای نمونه‌ها.
	}
	\label{q1:gradcam}
	
\end{figure}

همان‌طور که در سؤال خواسته شد، 
\lr{Grad-CAM}
در نوت‌بوک 
\lr{q2\_gradcam}
پیاده‌سازی شد. در تصویر 
\ref{q1:gradcam}
می‌توان دید که مدل روی چه نواحی‌ای تمرکز کرده است. برای مثال، در تصویر مربوط به سمندر، مدل روی نواحی‌ای که حیوان در آن‌ها قرار دارد تمرکز بیشتری نشان می‌دهد. 

البته می‌توان دید که مدل کاملاً دقیق نیست و برای مثال در تصاویر مربوط به سطل و ماسک تنها روی بخش کوچکی از تصویر تمرکز دارد. 

با پیاده‌سازی 
\lr{Guided Grad-CAM}
و
\lr{Guided Backpropagation}
و مقایسه‌ی خروجی آن‌ها که در شکل 
\ref{q1:guidedgradcam}
آمده است، می‌توان دید که خروجی‌ها تقریباً شبیه لبه‌ها و ساختارهای محلی تصویر هستند که مطابق انتظار است. 

یک تفاوت جالب که در تصویر عروس دریایی دیده می‌شود این است که در 
\lr{Guided Backpropagation}
خروجی دارای مؤلفه‌های آبی و قرمز بیشتری است که در تصویر نهایی به صورت رنگ‌های متمایل به بنفش دیده می‌شوند؛ ولی در 
\lr{Guided Grad-CAM}
از آنجا که مدل تشخیص می‌دهد زمینه‌ی آبی نقش زیادی در نواحی مؤثر برای تصمیم‌گیری ندارد، بخش‌های سبز و قرمز بیشتر برجسته می‌شوند. 

تفاوت دیگر که در خروجی‌های مقاله نیز مشاهده می‌شود این است که 
\lr{Guided Backpropagation}
لبه‌ها را بهتر مشخص می‌کند، ولی 
\lr{Guided Grad-CAM}
فقط ساختارهایی را برجسته می‌کند که در انتخاب کلاس نیز مؤثر هستند. 

\begin{figure}[H]
	
	\centering
	\includegraphics[width=\textwidth]{images/q1_guidedgradcam.png}
	\caption{
		خروجی 
		\lr{Guided Grad-CAM}
		برای نمونه‌ها.
	}
	\label{q1:guidedgradcam}
	
\end{figure}

\subsection{\lr{FinerGrad-CAM}}

\subsubsection{پرسش و پاسخ}

\begin{enumerate}
	
	\item 
	زیرا 
	\lr{Grad-CAM}
	نواحی‌ای را نشان می‌دهد که باعث تقویت کلاس شده‌اند، اما این نواحی لزوماً متمایزکننده\LTRfootnote{Discriminative} نیستند. 
	
	مثالی که در مقاله مطرح می‌شود مربوط به دسته‌بندی پرندگان است؛ ممکن است دو کلاس مختلف هر دو به بخش آبی بدن پرنده وزن بالایی بدهند. در نتیجه، این ویژگی اگرچه برای هر دو کلاس مهم است، اما الزاماً عامل تمایز میان آن‌ها نیست.
	
	\item 
	ادعای مطرح‌شده در مقاله این است که مغز هنگام تشخیص یک نمونه در میان گروه‌های مختلف، به دنبال ویژگی‌های متمایزکننده می‌گردد. 
	
	\lr{FinerGrad-CAM}
	نیز با نشان دادن نواحی‌ای که احتمال کلاس هدف را بدون تقویت سایر کلاس‌های مشابه افزایش می‌دهند، تلاش می‌کند ویژگی‌های یکتای متمایزکننده را مشخص کند. 
	
	به جای استفاده‌ی مستقیم از گرادیان کلاس هدف، از گرادیان‌های کنتراستی استفاده می‌شود تا ویژگی‌هایی که بین کلاس هدف و کلاس‌های مشابه مشترک هستند وزن کمتری بگیرند.
	
	\item 
	هرچه مقدار 
	$\gamma$
	بیشتر باشد، 
	\lr{FinerGrad-CAM}
	نواحی کوچک‌تر و متمایزکننده‌تری را نمایش می‌دهد؛ یعنی نواحی‌ای که فقط احتمال کلاس هدف را افزایش می‌دهند. 
	
	اما در مقابل، بخشی از اطلاعات کلی مورد استفاده‌ی مدل برای انتخاب کلاس از دست می‌رود.
	
	برای مثال، در کاربردهای زیست‌شناسی اگر هدف یافتن مهم‌ترین نواحی مؤثر در پیش‌بینی کلاس باشد، مقدار گاما باید کوچک انتخاب شود. ولی اگر هدف یافتن ویژگی‌هایی باشد که باعث ترجیح یک کلاس نسبت به کلاس‌های مشابه شده‌اند، استفاده از مقادیر بزرگ‌تر گاما مناسب‌تر است.
	
	\item 
	در مثالی که مقاله مطرح می‌کند، در یک مجموعه‌داده‌ی تصویر-توضیح بررسی می‌شود که آیا مدل به بخش‌هایی از تصویر که در توضیح ذکر شده‌اند توجه می‌کند یا خیر. 
	
	این موضوع باعث می‌شود فارغ از دقت نهایی مدل، اعتمادپذیری تصمیم‌های آن افزایش یابد و رفتار مدل پایدارتر باشد. 
	
	اما اگر مدل بدون توجه به توضیحات به پاسخ صحیح برسد، احتمال بیش‌برازش روی یک ویژگی خاص افزایش می‌یابد و سامانه احتمالاً از پایداری و مقاومت کمتری برخوردار خواهد بود.
	
\end{enumerate}

\subsubsection{پیاده‌سازی}

\begin{figure}[H]
	
	\centering
	\includegraphics[width=\textwidth]{images/q1_finer3.png}
	\caption{
		نمونه تصویر گروه ۳ که برای 
		\lr{d}
		تعیین شده است. این گروه مربوط به «والیبال» است.
	}
	\label{q2:fgcam}
	
\end{figure}

\begin{figure}[H]
	
	\centering
	\includegraphics[width=\textwidth]{images/q1_finerresults.png}
	\caption{
		خروجی 
		\lr{FinerGrad-CAM}
		برای نمونه‌ها.
	}
	\label{q2:fgcam_res}
	
\end{figure}

همان‌طور که در تصویر 
\ref{q2:fgcam_res}
مشاهده می‌شود، با افزایش گاما، نواحی مشخص‌شده کوچک‌تر می‌شوند و ویژگی‌های شاخص‌تر تصویر باقی می‌مانند. 

البته این مجموعه‌داده و این شیوه‌ی آزمایش، داده‌ی مناسبی برای نمایش کامل توانایی
\lr{FinerGrad-CAM}
نیست؛ زیرا کلاس‌ها شباهت زیادی به یکدیگر ندارند و برای مشاهده‌ی تغییرات محسوس، باید از مقادیر بسیار بزرگ گاما استفاده شود. 

(برای مثال در اینجا از مقادیر ۱، ۵ و ۱۰۰ استفاده شد، در حالی که مقدار متعارف گاما معمولاً بین صفر و یک است.)

در تصویر 
\ref{q2:fgcam_res}
خروجی مربوط به 
$\gamma = 0$
نیز رسم شده است که عملاً همان خروجی 
\lr{Grad-CAM}
است. 

نکته‌ی قابل مشاهده‌ی دیگر در تصویر مربوط به کوالا است. می‌توان دید که با افزایش زیاد گاما، ناحیه‌ی قبلی دیگر برجسته نمی‌شود و مدل روی سایر نواحی متمایزکننده‌ای که در خروجی اولیه چندان مشخص نبودند نیز تمرکز می‌کند. این موضوع ایده‌ی اصلی 
\lr{FinerGrad-CAM}
را به‌خوبی نشان می‌دهد.

\subsection{مقایسه}

در تصویر 
\ref{q2:fgcam_res}
خروجی سطر مربوط به 
$\gamma = 0$
برابر خروجی 
\lr{Grad-CAM}
و خروجی سطر مربوط به 
$\gamma = 1$
برابر با خروجی استاندارد 
\lr{FinerGrad-CAM}
است. 

(سایر مقادیر گاما خارج از بازه‌ی متعارف هستند و صرفاً برای ایجاد تغییرات قابل مشاهده انتخاب شده‌اند.)

می‌توان دید که در 
\lr{FinerGrad-CAM}
در بعضی تصاویر، مانند تصویر سطل، ناحیه‌ی کوچک‌تری مشخص شده است. این موضوع نشان می‌دهد که بخش‌هایی از ناحیه‌ی اولیه احتمالاً چندان متمایزکننده نبوده‌اند و احتمال چند کلاس مختلف را هم‌زمان افزایش می‌داده‌اند. 

در مقابل، در برخی تصاویر مانند کوالا یا ماسک، ناحیه‌ی برجسته‌شده بزرگ‌تر شده است. وجود این ناحیه‌ی جدید که در خروجی قبلی مشخص نبود، نشان می‌دهد بخشی از ناحیه‌ی اولیه میان چند کلاس مشترک بوده است. بنابراین با استفاده از گرادیان‌های کنتراستی و حذف فعال‌سازی‌های مشترک بین کلاس‌ها، می‌توان نواحی واقعاً متمایزکننده را بهتر مشخص کرد. 

همچنین در تصویر سمندر (و تا حدی عروس دریایی) مشاهده می‌شود که افزایش گاما تأثیر زیادی بر خروجی ندارد، مگر در مقادیر بسیار بزرگ. این موضوع نشان می‌دهد که مدل از ابتدا روی ویژگی‌های نسبتاً یکتای سمندر تمرکز کرده است و ویژگی مشترک قابل‌توجهی با کلاس مرجع انتخاب‌شده وجود ندارد.